\section{Methodology}
\label{sec:method}

\subsection{Dataset Construction}

The dataset pairs 261 expert-written .feature files from three open-source projects with 261 user stories. The stories were not collected from a backlog. They were reverse-engineered from the feature files: GPT-5.4-mini (temperature 0) received the text of each feature file and was asked to write one Connextra-style story (``As a [role], I want [feature], so that [benefit]'') with as many acceptance criteria as needed to cover all of its scenarios. Only the first 15,000 characters of each file were passed to the model. This truncated 84 of the 261 files, all of them from Apache Fineract (84 of its 100 files, whose median length is about 30,000 characters). The full feature file is always used as the reference. Section~\ref{sec:threats} discusses the leakage risk this construction creates, and Section~\ref{sec:res-b1} measures it.

\begin{table}[t]
\centering
\caption{Dataset composition. The three exemplars of the original prompt (one per project) are excluded from the primary analysis.}
\label{tab:dataset}
\begin{tabular}{llrr}
\hline
Project & Domain & Feature files & Analysed\\
\hline
Sylius & E-commerce & 100 & 99\\
Apache Fineract & Finance / banking & 100 & 99\\
Diaspora & Social network & 61 & 60\\
\hline
Total & & 261 & 258\\
\hline
\end{tabular}
\end{table}

The unit of analysis is a feature file, not a scenario. With the behave 1.3.3 parser, the 244 parseable files among the 258 analysed references contain 2,554 expanded scenarios, and the 255 parseable generated files contain 1,913. The reference files were not screened with a parser before inclusion: 14 of the 261 reference files fail to parse with behave 1.3.3 (13 with behave 1.2.6). They remain in the dataset, because the metric of Section~\ref{sec:metric1} does not require a parse.

The projects were chosen for domain and style diversity. Sylius covers products, carts, channels, shipping and taxes. Fineract contains data-heavy loan workflows with dates, balances, repayment schedules and data tables. Diaspora covers social-network behaviour driven through the browser. The benchmark is a deliberately heterogeneous sample of mature public English-language repositories, not a representative sample of BDD practice.

\subsection{Model and Prompt Configuration}

All generations use the snapshot gpt-5.4-mini-2026-03-17 through the OpenAI Chat Completions API with temperature 0. The requests set no maximum-token cap, which deviates from the research proposal's planned cap of 1,500 output tokens. The original condition (C0) is same-project one-shot prompting: each project has its own system message and one example pair (user story and expert Gherkin) from the same project (Listing~\ref{lst:prompt} in Appendix~\ref{app:details}). The three example files (ids 2, 102 and 202) are also members of the dataset, and their own generated files score 1.000, 0.997 and 1.000 on the metric below. The primary analysis therefore excludes them ($N = 258$); results for all 261 files are reported alongside.

\textbf{Prompt ablation (added in revision).} Two further conditions test how much the results depend on the example. C1 is zero-shot: the same system message and wording, with the example block and the sentence referring to it removed. C2 is one-shot with a different exemplar per project (ids 86, 119 and 203), drawn with a fixed random seed (2026) from the remaining files of each project before any C2 output existed. Model, temperature, system messages and output clean-up are identical in all conditions. The ablation is analysed on the 255 files that are not an exemplar in any condition. Its analysis plan was written before the C1 and C2 runs and is included in the replication package with its SHA-256 hash. It specifies paired Wilcoxon signed-rank tests on per-file cosine and exact McNemar tests on parser outcomes, each against C0 and Holm-corrected. The C1 and C2 runs took place on 16 October 2026. Because the environment used for this analysis could not download the encoder, skeleton cosine for all three conditions was computed with an ONNX export of all-MiniLM-L6-v2 using the same tokenisation, 256-token truncation, mean pooling and normalisation; it reproduces the original C0 scores to within $4 \times 10^{-7}$, and a rerun with sentence-transformers (model revision \texttt{1110a24}) agrees with it to within $3 \times 10^{-7}$ in all three conditions.

\subsection{Generation and Recording Pipeline}

The pipeline loads the stories, builds the prompt for the story's project and calls the API. It removes a Markdown code fence only when the model wraps its output in one; no generated text is edited by hand. Failed requests are retried for the same story. In C0 all 261 requests returned an output, so no file was dropped. Each call is appended to a log with its timestamp and token counts. In the C0 log, the model field is the requested model name written by the script, not the name returned by the API, and the logged cost was computed with outdated prices. Recomputed from the logged token counts (328,257 prompt and 380,494 completion tokens) at the Standard list price of \$0.75 and \$4.50 per million input and output tokens, without cache discounts, the estimated C0 cost is approximately \$1.96 (an estimate from the logged tokens, not an invoice). The C1 and C2 runs log the model name returned by the API.

\subsection{Metric 1: Skeleton Cosine Similarity}
\label{sec:metric1}

Skeleton cosine similarity compares a generated file with its reference after reducing both to a skeleton. The skeleton keeps every line that begins with a Gherkin keyword (Feature, Background, Scenario, Given, When, Then, And, But), including the keyword itself. It removes tags, replaces every double-quoted string with \texttt{<PARAM>} and every number with \texttt{<NUM>}, and drops all other lines, so data tables, doc strings and Examples tables disappear (Table~\ref{tab:skeleton} in Appendix~\ref{app:details}). Each skeleton is encoded as a whole with all-MiniLM-L6-v2~\cite{reimers2019}, which mean-pools token embeddings and truncates its input at 256 word-piece tokens. With $g$ and $r$ the embeddings of the generated and reference skeletons,
\begin{equation}
\mathrm{SkelCos}(g, r) = \frac{g \cdot r}{\lVert g \rVert \lVert r \rVert}.
\end{equation}

Two properties of this design limit what the metric can show. First, masking makes the score blind to changed numbers and quoted values by construction. Second, truncation means that the score often reflects only the start of a file: 174 of the 258 reference skeletons (67.4\%) and 166 of the 258 generated skeletons (64.3\%) exceed 256 tokens, including 94 of the 99 Fineract references (median 3,506 tokens). Section~\ref{sec:res-defects} therefore checks the metric against defect evidence instead of treating it as a measure of business correctness.

The threshold of 0.85 was fixed in the research proposal before the full run. It was not calibrated against human judgement. No prior study reports skeleton cosine with this encoder for Gherkin, and the proposal's pilot (30 stories, ids 1--10, 101--110 and 201--210, mean 0.859) used stories that are also part of the evaluated set. We therefore treat 0.85 as an operational reference point and report how the conclusions change between 0.80 and 0.90 (Section~\ref{sec:res-b3}).

\subsection{Metric 2: Parser Compliance}

Parser compliance is the share of generated files that the behave Gherkin parser accepts. For a generated set $S$,
\begin{equation}
\mathrm{ESR} = \frac{|\{s \in S \mid \text{parse}(s) = \text{PASS}\}|}{|S|}.
\end{equation}
We keep the proposal's name, executable syntax rate (ESR), but the check is narrower than the name suggests. Each file is passed to behave's parser (behave 1.3.3); no step definitions exist, so no step is matched or executed. A PASS means only that the file is valid Gherkin. It does not mean that the scenarios are complete, executable against the application or correct. The submitted version ran \mbox{behave -{}-dry-run} for this check. Without a steps directory, behave stops with a configuration error before parsing the file, and the original script counted that outcome as a PASS. All parser results in this version were therefore recomputed with the parser itself; behave 1.2.6 and 1.3.3 agree on every generated file.

\subsection{Defect Audit and Structural Checks (added in revision)}
\label{sec:method-audit}

Because neither metric checks business intent, we audited high-scoring outputs for defects. Of the 258 analysed files, 136 have a skeleton cosine of at least 0.85. From these we drew 20 files per project with a fixed seed (2026), giving 60 files. Each file was labelled against its user story and acceptance criteria; the reference was consulted only for project conventions. For each criterion the labeller recorded two facets. The business-logic facet covers wrong or missing numbers, operators, negations and conditions, wrong expected outcomes, untested criteria and contradicting additions. The temporal/causal facet covers missing preconditions, wrong triggering actions, missing or non-causal outcomes and impossible state orderings. Each defect received a severity from 0 (none) to 3 (the test would pass on a system that violates the criterion), and a file counts as defective if it has at least one defect of severity 2 or higher. Line-level evidence was recorded for each defect.

The labels were drafted by an AI coding assistant (OpenAI Codex) and then reviewed by a second author, who retained all of them unchanged. That review was recorded as one confirmation covering all 60 files, not as a decision with reasons per file. To estimate reliability, the first author, who reports not having seen the labels or their aggregate rates before rating, labelled a subset of 15 files fixed in advance (6 Sylius, 3 Fineract, 6 Diaspora) with the same guide. We report Cohen's $\kappa$~\cite{cohen1960} between these independent labels and the reviewed labels. This is agreement between one independent human rater and AI-drafted, human-reviewed labels, not agreement between two independent human raters.

Separately, an automated check parses each file and flags four keyword-order violations per scenario: no Then step, a Given after a When or Then, a Then before any When, and a non-Given step in the Background. These flags encode the checker's conventions and do not by themselves establish a business defect.

Finally, a perturbation test applied eight defect-seeding operators to 120 scenarios extracted from the reference files and compared each mutant with its original using the same metric (Appendix~\ref{app:perturb}).

\subsection{Statistical Analysis Plan}

The confirmatory tests use $\alpha = 0.05$ and follow standard guidance on matching tests to data properties~\cite{arcuri2011}.

\textbf{RQ1: one-sample Wilcoxon signed-rank test.} The test is applied to $x_i - 0.85$, one-sided, with zero differences dropped. $H_0$: the median skeleton cosine is at most 0.85; $H_1$: it exceeds 0.85. The proposal phrased these hypotheses in terms of the mean; we state them for the median, which the signed-rank test addresses under the assumption that the differences are symmetric. As effect size we report $\delta = P(x > 0.85) - P(x < 0.85)$, a one-sample analogue of Cliff's $\delta$ against the threshold. Quartiles use linear interpolation and the SD uses $n - 1$. A non-significant result is read as insufficient evidence for performance above the threshold, not as evidence of poor performance.

\textbf{RQ2: exact binomial test.} Each file has a binary parser outcome. We test $H_0$: $p_{\mathrm{ESR}} \le 0.80$ against $H_1$: $p_{\mathrm{ESR}} > 0.80$~\cite{agresti2002} and report a 95\% Clopper--Pearson interval~\cite{clopper1934}.

\textbf{Analyses added in revision.} The following analyses were planned after the reviews and are exploratory. Per-project differences are tested with a Kruskal--Wallis test~\cite{kruskal1952}, reported with $\eta^2_H = (H - k + 1)/(n - k)$ and followed by pairwise two-sided Mann--Whitney tests with Holm correction~\cite{holm1979}. Threshold sensitivity repeats the RQ1 test for thresholds from 0.80 to 0.90. Lexical leakage is measured as the share of distinct word trigrams of the reference file that also occur in the user story, and is related to skeleton cosine with Spearman's $\rho$.

\subsection{Reproducibility and Denominators}

The replication package contains the dataset, prompts, scripts, generated outputs, API logs and all revision analyses at \url{https://github.com/simonhoang611/RT-SWT-004-Group4}. The primary analysis uses $N = 258$ files, and no generated file was excluded because of its content; only the three exemplars are left out of the primary analysis.
